\section{The model patch and the height equation}\label{sec:model}

The first task is to obstruct scalar height functions on a whole
closed square. We fix a curvature profile with a negative center
and a positive exterior, then specify the equation and inequalities
that the final geometric argument will force on one such square.

Fix a smooth even function \(h:\R\to\R\) such that
\begin{equation}\label{eq:h}
\begin{gathered}
h=1 \quad\text{on }[-1/2,1/2],\qquad h'\le0\quad\text{on }[0,\infty),\\
h>0\quad\text{on }(-1,1),\qquad h<0\quad\text{for }\abs y>1,\\
h(1)=0,\qquad h'(1)<0.
\end{gathered}
\end{equation}
For example, let \(\rho\) be a nonnegative smooth function supported
in \([1/2,2]\), positive on \((1/2,3/2)\), and put
\[
h(y)=1-\frac{\int_0^{\abs y}\rho(v)\dd v}
                  {\int_0^1\rho(v)\dd v}.
\]
This function has all the required properties.

Write \(S=[-3,3]^2\), with coordinates \(x,y\), and fix a positive
constant \(\kappa\). Let \(g_0\) be a smooth positive metric on a
neighborhood of \(S\) with curvature
\begin{equation}\label{eq:model-curvature}
K_{g_0}(x,y)=\kappa\bigl(x^2-h(y)\bigr).
\end{equation}
We will construct such background metrics in \Cref{sec:assembly}.
Perturbations are smooth symmetric tensors supported in
\[
\mathcal P=[-1/10,1/10]^2.
\]
Let \(\mathcal U\) be the open set of these perturbations for which
the resulting metric \(g\) is positive and
\begin{equation}\label{eq:negative-box}
K_g<-\kappa/2\qquad\text{on }[-1/5,1/5]^2.
\end{equation}
We regard \(\mathcal U\) as a set of metrics with its \(C^\infty\)
topology. It contains \(g_0\): on this smaller square,
\(h=1\) and \(x^2\le1/25\). Outside \(\mathcal P\), every metric
in \(\mathcal U\) still has curvature \eqref{eq:model-curvature}.

For a function \(z\), write
\begin{equation}\label{eq:height-notation}
H_{ij}=z_{ij}-\Gamma^r_{ij}z_r,\qquad
E=\det(g)(1-\abs{dz}_g^2).
\end{equation}
Here \(K=K_g\), \(\Gamma^r_{ij}\) are the Christoffel symbols of
\(g\), and repeated indices are summed. Wherever \(H_{yy}\ne0\),
put \(q=H_{xy}/H_{yy}\). An \emph{admissible height}
for \(g\) is a function \(z\), smooth on a neighborhood of the closed square \(S\), satisfying throughout~\(S\)
\begin{equation}\label{eq:admissible}
H_{xx}H_{yy}-H_{xy}^2=KE,\qquad
E>0,\qquad H_{yy}\ne0,\qquad \abs q\le1/100.
\end{equation}
The scalar equation here is the Darboux equation; compare
\cite[Equation~(1) and Appendix]{Khuri2009}. We recall the necessary geometric
identity directly.

\begin{lemma}[Height equation]\label{lem:height}
Let \(F\) be a smooth isometric immersion of a surface into \(\R^3\),
let \(e\) be a fixed ambient unit vector, and set \(z=\langle F,e\rangle\).
Then
\[
\det(\nabla_g^2z)=K\det(g)(1-\abs{dz}_g^2).
\]
If \(e\) is normal to the immersed surface at a point \(b\), then
\(dz(b)=0\) and \(\nabla_g^2z(b)=\II(b)\), with the choice of unit
normal equal to \(e\) at \(b\).
\end{lemma}

\begin{proof}
With \(n\) the chosen unit normal and \(\nu=\langle n,e\rangle\),
the Gauss formula gives \(\nabla_g^2z=\nu\II\). The tangential and
normal decomposition of \(e\) gives
\(\abs{dz}_g^2+\nu^2=1\).
The Gauss equation \(\det\II=K\det g\) now proves the identity,
and the final statement follows from \(\nu(b)=1\).
\end{proof}

Only this necessary equation for a height is used. We do not need
the converse reconstruction of an immersion from a Darboux solution.

\begin{proposition}[Patch obstruction]\label{prop:patch}
The metrics in \(\mathcal U\) admitting no admissible height form
a residual subset of \(\mathcal U\). In particular, for every
\(C^\infty\) neighborhood of \(g_0\) in \(\mathcal U\), there is
a metric in that neighborhood with no admissible height.
\end{proposition}

The whole-square requirement is essential. Every curvature zero
in \(S\) of a metric in \(\mathcal U\) is a clean change of sign. Indeed,
\(K<0\) on \(\mathcal P\); outside it,
\(K_x=2\kappa x\), and a zero with \(x=0\) has
\(y=\pm1\), where \(K_y\ne0\). Thus the pointwise local-existence
results recalled in the introduction are compatible with this
proposition: they do not supply one height satisfying
\eqref{eq:admissible} throughout \(S\).

We prove \Cref{prop:patch} in \Cref{sec:pulse}, after establishing estimates on
the two caps. \Cref{fig:geometry} shows the model geometry.

\begin{figure}[H]
\centering
\begin{tikzpicture}[font=\small,>=Latex]
  \begin{scope}[scale=1.12]
    \fill[blue!5] (-1.7,-1.7) rectangle (1.7,1.7);
    \fill[orange!18]
      (0,1.35) .. controls (.80,1.35) and (1.02,.95) .. (1.02,.60)
      -- (1.02,-.60)
      .. controls (1.02,-.95) and (.80,-1.35) .. (0,-1.35)
      .. controls (-.80,-1.35) and (-1.02,-.95) .. (-1.02,-.60)
      -- (-1.02,.60)
      .. controls (-1.02,.95) and (-.80,1.35) .. (0,1.35);
    \draw[thick]
      (0,1.35) .. controls (.80,1.35) and (1.02,.95) .. (1.02,.60)
      -- (1.02,-.60)
      .. controls (1.02,-.95) and (.80,-1.35) .. (0,-1.35)
      .. controls (-.80,-1.35) and (-1.02,-.95) .. (-1.02,-.60)
      -- (-1.02,.60)
      .. controls (-1.02,.95) and (-.80,1.35) .. (0,1.35);
    \draw[gray] (-1.7,-1.7) rectangle (1.7,1.7);
    \draw[->,gray] (-1.87,0)--(1.91,0) node[right,black] {\(x\)};
    \draw[->,gray] (0,-1.87)--(0,1.91) node[above,black] {\(y\)};
    \fill[white] (-.25,-.25) rectangle (.25,.25);
    \draw[thick] (-.25,-.25) rectangle (.25,.25);
    \node at (.02,.01) {\(\mathcal P\)};
    \node at (.0,.76) {\(K<0\)};
    \node at (1.28,1.42) {\(K>0\)};
    \node[fill=white,inner sep=1pt] at (.90,-1.15) {\(K=0\)};
  \end{scope}
\end{tikzpicture}
\caption{The model patch, schematically and not to scale.
The curvature zero set lies outside the perturbation box
\(\mathcal P\). The positive exterior permits estimates on lower and
upper caps whose free artificial edges may cut through the negative
region. The lower-cap integration region is shown in \Cref{fig:cap}.}
\label{fig:geometry}
\end{figure}

